\section{Conclusão}\label{sec:conclusao}

Este trabalho propôs e avaliou o \emph{AIMD CUBIC}, um
controlador adaptativo de taxa de bits para o servidor de
\emph{streaming} de jogos em nuvem Sunshine. O controlador
estende um esquema AIMD com duas primitivas do algoritmo
CUBIC~\cite{cubic}, a curva cúbica de recuperação após
reduções e o mecanismo de \emph{fast convergence}, e opera
exclusivamente a partir do RTT mantido pela biblioteca ENet
no canal de controle, sem qualquer cooperação adicional do
cliente. Avaliado contra um controlador de referência fixa e
uma variante AIMD de base ao longo de seis cenários de
degradação de rede, o controlador proposto reduz o pico de
RTT em cinco dos seis cenários, com a maior queda alcançando
43\% no cenário~A, e o tempo total acumulado em \emph{stall}
em cinco dos seis cenários, com queda de 61\% no cenário~E.
Esses ganhos são acompanhados de uma elevação sistemática do
RTT mediano, de 4--5~ms no AIMD de base para 12--23~ms no
AIMD CUBIC, configurando um trade-off explícito entre
desempenho na cauda da distribuição de latência e desempenho
no regime estacionário. Os dois controladores representam,
portanto, escolhas de prioridade complementares, não soluções
mutuamente dominantes.

\subsection{Achados que Transcendem o Sunshine}

A análise de sensibilidade ao design da curva (§\ref{sec:sensibilidade})
isolou um achado de alcance maior que a contribuição
imediata. Variantes mais fiéis à formulação original do CUBIC,
incluindo a interpretação literal do parágrafo §3.1 do
trabalho original, mostraram-se sistematicamente inferiores à
forma simplificada adotada neste trabalho. O diagnóstico é
estrutural: a curva cúbica do CUBIC foi formulada para
amostragem por~ACK, regime em que pequenos deltas se acumulam
ao longo de muitas iterações por segundo. Em controladores
amostrados a cadências reduzidas, como o aqui proposto em
1~Hz, a curva passa a aplicar o delta integral entre instantes
de decisão, e a fidelidade ao paper original deixa de ser
empiricamente desejável. A transposição de algoritmos de
controle de congestionamento entre regimes de cadência exige,
portanto, ajuste estrutural, não apenas reparametrização.

Dois outros achados merecem destaque. O primeiro é que o RTT
mantido server-side pela biblioteca ENet é, isoladamente,
sinal suficiente para um controlador útil em jogos em nuvem,
sem que sejam necessárias extensões de protocolo nem
sincronização de relógio entre cliente e servidor. O segundo é
que métricas de cauda (pico, percentil~99, eventos de
\emph{stall}) e métricas centrais (mediana, percentis
intermediários) devem ser reportadas conjuntamente para
caracterização correta de controladores adaptativos: nenhuma
das duas isoladamente é suficiente para distinguir prioridades
de projeto, e o desempenho subjetivo correlaciona de forma
não-trivial com ambas.

\subsection{Trabalhos Futuros}

Cinco direções concretas emergem como continuação natural
deste trabalho. A primeira é a introdução de uma memória de
capacidade efetiva entre sessões, motivada pela regressão
observada no cenário~F. Naquele cenário, a curva inicial de
recuperação fica apontada para o teto negociado de sessão,
inalcançável dada a constrição ativa desde $t=0$. Uma
estimativa de capacidade carregada da sessão anterior
permitiria que o controlador iniciasse com um teto
\texttt{last\_known\_good} mais próximo da capacidade real do
caminho, evitando a primeira cascata desperdiçada.

A segunda direção é a adoção de um sinal de detecção
sub-segundo, possivelmente a partir do tempo de envio do
quadro mais recente (\texttt{frame\_send\_ms}, já coletado na
instrumentação atual). A latência de reação do controlador
proposto é da ordem da histerese de decremento, e um sinal
adicional que se manifeste antes da inflação do RTT pode
encurtar o tempo entre o início do congestionamento e o corte
de taxa.

A terceira direção é a introdução de um gatilho preditivo
baseado em gradiente, na linha do controlador GCC para
WebRTC~\cite{gcc}, capaz de antecipar a transição da fase de
não-congestionamento para a fase de enfileiramento antes que
o RTT cruze o limiar.

A quarta direção é a condução de uma avaliação subjetiva
\emph{crowdsourced}, segundo o protocolo de Hong et
al.~\cite{hong2015}, em substituição às avaliações de
operador único conduzidas neste trabalho. Esse formato
permite quantificação estatística de Mean Opinion Score sob
diferentes perfis de jogador, tipos de jogo, e condições de
rede.

A quinta direção é a comparação direta do AIMD CUBIC contra o
GCC~\cite{gcc} e o ABRAGame~\cite{abragame} em um mesmo
testbed, com a mesma instrumentação. Embora os trabalhos
correlatos tenham reportado seu desempenho em ambientes
próprios, a comparação cruzada em ambiente comum é a única
forma robusta de posicionar o presente trabalho contra o
estado da arte em métricas absolutas.

\subsection{Impacto}

O controlador proposto, juntamente com a instrumentação de
coleta de métricas que o sustenta, será submetido como
contribuição ao repositório \emph{upstream} do Sunshine sob a
licença GPL-3.0 do projeto. O testbed experimental, os
\emph{scripts} de geração de cenários, os dados de cada
sessão e os \emph{notebooks} de análise estão disponíveis para
reprodução. As restrições metodológicas adotadas, como a
execução única por cenário, a padronização da carga de
trabalho em um único título e a avaliação subjetiva por um
único operador, são documentadas ao longo do trabalho e
abrem o caminho natural para as direções de trabalho futuro
listadas acima.
